\IfFileExists{stacks-project.cls}{%
\documentclass{stacks-project}
}{%
\documentclass{amsart}
}

% For dealing with references we use the comment environment
\usepackage{verbatim}
\newenvironment{reference}{\comment}{\endcomment}
%\newenvironment{reference}{}{}
\newenvironment{slogan}{\comment}{\endcomment}
\newenvironment{history}{\comment}{\endcomment}

% For commutative diagrams we use Xy-pic
\usepackage[all]{xy}

% We use 2cell for 2-commutative diagrams.
\xyoption{2cell}
\UseAllTwocells

% We use multicol for the list of chapters between chapters
\usepackage{multicol}

% This is generally recommended for better output
\usepackage{lmodern}
\usepackage[T1]{fontenc}

% For cross-file-references
\usepackage{xr-hyper}

% Package for hypertext links:
\usepackage{hyperref}

% For any local file, say "hello.tex" you want to link to please
% use \externaldocument[hello-]{hello}
\externaldocument[introduction-]{introduction}
\externaldocument[conventions-]{conventions}
\externaldocument[sets-]{sets}
\externaldocument[categories-]{categories}
\externaldocument[topology-]{topology}
\externaldocument[sheaves-]{sheaves}
\externaldocument[sites-]{sites}
\externaldocument[stacks-]{stacks}
\externaldocument[fields-]{fields}
\externaldocument[algebra-]{algebra}
\externaldocument[brauer-]{brauer}
\externaldocument[homology-]{homology}
\externaldocument[derived-]{derived}
\externaldocument[simplicial-]{simplicial}
\externaldocument[more-algebra-]{more-algebra}
\externaldocument[smoothing-]{smoothing}
\externaldocument[modules-]{modules}
\externaldocument[sites-modules-]{sites-modules}
\externaldocument[injectives-]{injectives}
\externaldocument[cohomology-]{cohomology}
\externaldocument[sites-cohomology-]{sites-cohomology}
\externaldocument[dga-]{dga}
\externaldocument[dpa-]{dpa}
\externaldocument[sdga-]{sdga}
\externaldocument[hypercovering-]{hypercovering}
\externaldocument[schemes-]{schemes}
\externaldocument[constructions-]{constructions}
\externaldocument[properties-]{properties}
\externaldocument[morphisms-]{morphisms}
\externaldocument[coherent-]{coherent}
\externaldocument[divisors-]{divisors}
\externaldocument[limits-]{limits}
\externaldocument[varieties-]{varieties}
\externaldocument[topologies-]{topologies}
\externaldocument[descent-]{descent}
\externaldocument[perfect-]{perfect}
\externaldocument[more-morphisms-]{more-morphisms}
\externaldocument[flat-]{flat}
\externaldocument[groupoids-]{groupoids}
\externaldocument[more-groupoids-]{more-groupoids}
\externaldocument[etale-]{etale}
\externaldocument[chow-]{chow}
\externaldocument[intersection-]{intersection}
\externaldocument[pic-]{pic}
\externaldocument[weil-]{weil}
\externaldocument[adequate-]{adequate}
\externaldocument[dualizing-]{dualizing}
\externaldocument[duality-]{duality}
\externaldocument[discriminant-]{discriminant}
\externaldocument[derham-]{derham}
\externaldocument[local-cohomology-]{local-cohomology}
\externaldocument[algebraization-]{algebraization}
\externaldocument[curves-]{curves}
\externaldocument[resolve-]{resolve}
\externaldocument[models-]{models}
\externaldocument[functors-]{functors}
\externaldocument[equiv-]{equiv}
\externaldocument[pione-]{pione}
\externaldocument[etale-cohomology-]{etale-cohomology}
\externaldocument[proetale-]{proetale}
\externaldocument[relative-cycles-]{relative-cycles}
\externaldocument[more-etale-]{more-etale}
\externaldocument[trace-]{trace}
\externaldocument[crystalline-]{crystalline}
\externaldocument[spaces-]{spaces}
\externaldocument[spaces-properties-]{spaces-properties}
\externaldocument[spaces-morphisms-]{spaces-morphisms}
\externaldocument[decent-spaces-]{decent-spaces}
\externaldocument[spaces-cohomology-]{spaces-cohomology}
\externaldocument[spaces-limits-]{spaces-limits}
\externaldocument[spaces-divisors-]{spaces-divisors}
\externaldocument[spaces-over-fields-]{spaces-over-fields}
\externaldocument[spaces-topologies-]{spaces-topologies}
\externaldocument[spaces-descent-]{spaces-descent}
\externaldocument[spaces-perfect-]{spaces-perfect}
\externaldocument[spaces-more-morphisms-]{spaces-more-morphisms}
\externaldocument[spaces-flat-]{spaces-flat}
\externaldocument[spaces-groupoids-]{spaces-groupoids}
\externaldocument[spaces-more-groupoids-]{spaces-more-groupoids}
\externaldocument[bootstrap-]{bootstrap}
\externaldocument[spaces-pushouts-]{spaces-pushouts}
\externaldocument[spaces-chow-]{spaces-chow}
\externaldocument[groupoids-quotients-]{groupoids-quotients}
\externaldocument[spaces-more-cohomology-]{spaces-more-cohomology}
\externaldocument[spaces-simplicial-]{spaces-simplicial}
\externaldocument[spaces-duality-]{spaces-duality}
\externaldocument[formal-spaces-]{formal-spaces}
\externaldocument[restricted-]{restricted}
\externaldocument[spaces-resolve-]{spaces-resolve}
\externaldocument[formal-defos-]{formal-defos}
\externaldocument[defos-]{defos}
\externaldocument[cotangent-]{cotangent}
\externaldocument[examples-defos-]{examples-defos}
\externaldocument[algebraic-]{algebraic}
\externaldocument[examples-stacks-]{examples-stacks}
\externaldocument[stacks-sheaves-]{stacks-sheaves}
\externaldocument[criteria-]{criteria}
\externaldocument[artin-]{artin}
\externaldocument[quot-]{quot}
\externaldocument[stacks-properties-]{stacks-properties}
\externaldocument[stacks-morphisms-]{stacks-morphisms}
\externaldocument[stacks-limits-]{stacks-limits}
\externaldocument[stacks-cohomology-]{stacks-cohomology}
\externaldocument[stacks-perfect-]{stacks-perfect}
\externaldocument[stacks-introduction-]{stacks-introduction}
\externaldocument[stacks-more-morphisms-]{stacks-more-morphisms}
\externaldocument[stacks-geometry-]{stacks-geometry}
\externaldocument[moduli-]{moduli}
\externaldocument[moduli-curves-]{moduli-curves}
\externaldocument[examples-]{examples}
\externaldocument[exercises-]{exercises}
\externaldocument[guide-]{guide}
\externaldocument[desirables-]{desirables}
\externaldocument[coding-]{coding}
\externaldocument[obsolete-]{obsolete}
\externaldocument[fdl-]{fdl}
\externaldocument[index-]{index}

% Theorem environments.
%
\theoremstyle{plain}
\newtheorem{theorem}[subsection]{Theorem}
\newtheorem{proposition}[subsection]{Proposition}
\newtheorem{lemma}[subsection]{Lemma}

\theoremstyle{definition}
\newtheorem{definition}[subsection]{Definition}
\newtheorem{example}[subsection]{Example}
\newtheorem{exercise}[subsection]{Exercise}
\newtheorem{situation}[subsection]{Situation}

\theoremstyle{remark}
\newtheorem{remark}[subsection]{Remark}
\newtheorem{remarks}[subsection]{Remarks}

\numberwithin{equation}{subsection}

% Macros
%
\def\lim{\mathop{\mathrm{lim}}\nolimits}
\def\colim{\mathop{\mathrm{colim}}\nolimits}
\def\Spec{\mathop{\mathrm{Spec}}}
\def\Hom{\mathop{\mathrm{Hom}}\nolimits}
\def\Ext{\mathop{\mathrm{Ext}}\nolimits}
\def\SheafHom{\mathop{\mathcal{H}\!\mathit{om}}\nolimits}
\def\SheafExt{\mathop{\mathcal{E}\!\mathit{xt}}\nolimits}
\def\Sch{\mathit{Sch}}
\def\Mor{\mathop{\mathrm{Mor}}\nolimits}
\def\Ob{\mathop{\mathrm{Ob}}\nolimits}
\def\Sh{\mathop{\mathit{Sh}}\nolimits}
\def\NL{\mathop{N\!L}\nolimits}
\def\CH{\mathop{\mathrm{CH}}\nolimits}
\def\proetale{{pro\text{-}\acute{e}tale}}
\def\etale{{\acute{e}tale}}
\def\QCoh{\mathit{QCoh}}
\def\Ker{\mathop{\mathrm{Ker}}}
\def\Im{\mathop{\mathrm{Im}}}
\def\Coker{\mathop{\mathrm{Coker}}}
\def\Coim{\mathop{\mathrm{Coim}}}

% Boxtimes
%
\DeclareMathSymbol{\boxtimes}{\mathbin}{AMSa}{"02}

%
% Macros for moduli stacks/spaces
%
\def\QCohstack{\mathcal{QC}\!\mathit{oh}}
\def\Cohstack{\mathcal{C}\!\mathit{oh}}
\def\Spacesstack{\mathcal{S}\!\mathit{paces}}
\def\Quotfunctor{\mathrm{Quot}}
\def\Hilbfunctor{\mathrm{Hilb}}
\def\Curvesstack{\mathcal{C}\!\mathit{urves}}
\def\Polarizedstack{\mathcal{P}\!\mathit{olarized}}
\def\Complexesstack{\mathcal{C}\!\mathit{omplexes}}
% \Pic is the operator that assigns to X its picard group, usage \Pic(X)
% \Picardstack_{X/B} denotes the Picard stack of X over B
% \Picardfunctor_{X/B} denotes the Picard functor of X over B
\def\Pic{\mathop{\mathrm{Pic}}\nolimits}
\def\Picardstack{\mathcal{P}\!\mathit{ic}}
\def\Picardfunctor{\mathrm{Pic}}
\def\Deformationcategory{\mathcal{D}\!\mathit{ef}}
\usepackage{cleveref}
\crefname{originallemma}{original lemma}{original lemmas}
\Crefname{originallemma}{Original lemma}{Original lemmas}
\makeatletter
\newlabel{lemma-pseudo-coherent-hocolim}{{0DJN}{source}{}{sourcepointer-1}{}}
\newlabel{lemma-pseudo-coherent-hocolim@cref}{{[originallemma][0][]0DJN}{[1][0][]source}{}{}{}}
\makeatother
\crefname{originallemma}{original lemma}{original lemmas}
\Crefname{originallemma}{Original lemma}{Original lemmas}
\makeatletter
\newlabel{cohomology-lemma-dual-perfect-complex}{{08DQ}{source}{}{sourcepointer-2}{}}
\newlabel{cohomology-lemma-dual-perfect-complex@cref}{{[originallemma][0][]08DQ}{[1][0][]source}{}{}{}}
\makeatother
\crefname{originallemma}{original lemma}{original lemmas}
\Crefname{originallemma}{Original lemma}{Original lemmas}
\makeatletter
\newlabel{lemma-perfect-relatively-perfect}{{0DI4}{source}{}{sourcepointer-3}{}}
\newlabel{lemma-perfect-relatively-perfect@cref}{{[originallemma][0][]0DI4}{[1][0][]source}{}{}{}}
\makeatother
\crefname{originallemma}{original lemma}{original lemmas}
\Crefname{originallemma}{Original lemma}{Original lemmas}
\makeatletter
\newlabel{lemma-derived-pushforward-rel-perfect}{{0DJT}{source}{}{sourcepointer-4}{}}
\newlabel{lemma-derived-pushforward-rel-perfect@cref}{{[originallemma][0][]0DJT}{[1][0][]source}{}{}{}}
\makeatother
\crefname{originallemma}{original lemma}{original lemmas}
\Crefname{originallemma}{Original lemma}{Original lemmas}
\makeatletter
\newlabel{lemma-tor-independence-and-tor-amplitude}{{0C0V}{source}{}{sourcepointer-5}{}}
\newlabel{lemma-tor-independence-and-tor-amplitude@cref}{{[originallemma][0][]0C0V}{[1][0][]source}{}{}{}}
\makeatother
\crefname{originallemma}{original lemma}{original lemmas}
\Crefname{originallemma}{Original lemma}{Original lemmas}
\makeatletter
\newlabel{more-algebra-lemma-lift-perfect-from-residue-field}{{0BCD}{source}{}{sourcepointer-6}{}}
\newlabel{more-algebra-lemma-lift-perfect-from-residue-field@cref}{{[originallemma][0][]0BCD}{[1][0][]source}{}{}{}}
\makeatother
\crefname{originallemma}{original lemma}{original lemmas}
\Crefname{originallemma}{Original lemma}{Original lemmas}
\makeatletter
\newlabel{cohomology-lemma-colim-and-lim-of-duals}{{0DJI}{source}{}{sourcepointer-7}{}}
\newlabel{cohomology-lemma-colim-and-lim-of-duals@cref}{{[originallemma][0][]0DJI}{[1][0][]source}{}{}{}}
\makeatother
\title{The derived sheaf Hom from a pseudo-coherent into an S-perfect complex is again a derived Hom from a pseudo-coherent complex}
\author{The Stacks Project Authors}
\date{stacks-project:perfect.tex:local-snapshot-20261012; GNU FDL 1.2 or later}
\begin{document}
\maketitle
\setcounter{subsection}{1}
\noindent\textit{Editorial note.} All delivered mathematical text is the author's own, copied from the pinned Stacks Project chapter file perfect.tex; no statement or proof step was rewritten, reordered or removed. The counted span is the author environment and its own proof environment. References inside the span are handed to the assembler as navigation stubs that carry the canonical tag of the cited tag only; whether each of them is load-bearing is recorded on the stub itself.\par\medskip
\noindent\textit{Allowed standard prerequisites (editorial).}\begin{itemize}
\item A pseudo-coherent complex on a quasi-compact quasi-separated scheme is a homotopy colimit of perfect complexes with isomorphisms on the truncations above a moving bound. invoked at proof lines 9251-9254. Produces the approximating perfect complexes K\_n from which the whole construction starts.
\item The dual of a perfect complex is perfect, and tensoring an S-perfect complex with a perfect complex stays S-perfect. invoked at proof lines 9254-9258. Makes the duals K\_n dual and their tensor products with E objects to which the relative perfection statement applies.
\item For a proper flat morphism of finite presentation the pushforward of an S-perfect complex is perfect and its formation commutes with arbitrary base change. invoked at proof lines 9258-9264. Produces the inverse system of perfect complexes M\_n on the base together with their base change compatibility.
\item Tor amplitude is independent of the base change, and a perfect complex whose fibre over a point has cohomology concentrated in high degrees has tor amplitude bounded below by that degree. invoked at proof lines 9275-9285. Gives the two sided tor amplitude bound for the cones C\_n, hence the bound for their duals.
\item The homotopy colimit of a system of perfect complexes with isomorphisms on truncations below a moving bound is pseudo-coherent. invoked at proof lines 9290-9295. Identifies the homotopy colimit of the duals L\_n with the required pseudo-coherent complex L.
\end{itemize}\medskip
\ifdefined\hypertarget\hypertarget{sourcepointer-1}{}\fi
\noindent\textit{Source reference.} Original lemma 0DJN: sources/stacks-perfect/pinned-original-source.tex, source lines 4814--4814. Invoked at proof lines 9251-9253 to obtain the approximating perfect complexes K\_n; the construction of the rest of the argument is inside the counted span.\par\medskip
\ifdefined\hypertarget\hypertarget{sourcepointer-2}{}\fi
\noindent\textit{Source reference.} Original lemma 08DQ: sources/stacks-chow/tags.txt, source lines 1--1. Cited at proof lines 9254-9256 to see that the duals K\_n dual are perfect; their use in the inverse system is inside the counted span.\par\medskip
\ifdefined\hypertarget\hypertarget{sourcepointer-3}{}\fi
\noindent\textit{Source reference.} Original lemma 0DI4: sources/stacks-perfect/pinned-original-source.tex, source lines 9007--9007. Used at proof lines 9256-9258 to see that K\_n dual tensor E is S-perfect; the application of the relative perfection statement that follows is inside the counted span.\par\medskip
\ifdefined\hypertarget\hypertarget{sourcepointer-4}{}\fi
\noindent\textit{Source reference.} Original lemma 0DJT: sources/stacks-perfect/pinned-original-source.tex, source lines 9188--9188. Applied at proof lines 9258-9264 to produce the system M\_n and its base change compatibility; the degree bookkeeping that follows is inside the counted span.\par\medskip
\ifdefined\hypertarget\hypertarget{sourcepointer-5}{}\fi
\noindent\textit{Source reference.} Original lemma 0C0V: sources/stacks-perfect/pinned-original-source.tex, source lines 5513--5513. Cited at proof lines 9273-9276 to see that the tor amplitude of E is unaffected by base change; the degree computation using it is inside the counted span.\par\medskip
\ifdefined\hypertarget\hypertarget{sourcepointer-6}{}\fi
\noindent\textit{Source reference.} Original lemma 0BCD: sources/stacks-chow/tags.txt, source lines 1--1. Invoked at proof lines 9283-9286 to bound the tor amplitude of the cone C\_n from its fibre; the use of that bound for the duals is inside the counted span.\par\medskip
\ifdefined\hypertarget\hypertarget{sourcepointer-7}{}\fi
\noindent\textit{Source reference.} Original lemma 0DJI: sources/stacks-chow/tags.txt, source lines 1--1. Cited at proof lines 9290-9295 for the identification of the homotopy colimit of the L\_n with the required complex; the pseudo-coherence conclusion drawn from it is inside the counted span.\par\medskip
\begin{lemma}
\label{lemma-compute-ext-rel-perfect}
Let $f : X \to S$ be a morphism of schemes. Let $E, K \in D(\mathcal{O}_X)$.
Assume
\begin{enumerate}
\item $S$ is quasi-compact and quasi-separated,
\item $f$ is proper, flat, and of finite presentation,
\item $E$ is $S$-perfect,
\item $K$ is pseudo-coherent.
\end{enumerate}
Then there exists a pseudo-coherent $L \in D(\mathcal{O}_S)$ such that
$$
Rf_*R\SheafHom(K, E) = R\SheafHom(L, \mathcal{O}_S)
$$
and the same is true after arbitrary base change: given
$$
\vcenter{
\xymatrix{
X' \ar[r]_{g'} \ar[d]_{f'} &
X \ar[d]^f \\
S' \ar[r]^g &
S
}
}
\quad\quad
\begin{matrix}
\text{cartesian, then we have } \\
Rf'_*R\SheafHom(L(g')^*K, L(g')^*E) \\
= R\SheafHom(Lg^*L, \mathcal{O}_{S'})
\end{matrix}
$$
\end{lemma}

\begin{proof}
Since $S$ is quasi-compact and quasi-separated, the same is true for $X$.
By Lemma \ref{lemma-pseudo-coherent-hocolim} we can write
$K = \text{hocolim} K_n$ with $K_n$ perfect and $K_n \to K$ inducing
an isomorphism on truncations $\tau_{\geq -n}$. Let $K_n^\vee$
be the dual perfect complex
(Cohomology, Lemma \ref{cohomology-lemma-dual-perfect-complex}).
We obtain an inverse system $\ldots \to K_3^\vee \to K_2^\vee \to K_1^\vee$
of perfect objects. By Lemma \ref{lemma-perfect-relatively-perfect}
we see that $K_n^\vee \otimes_{\mathcal{O}_X} E$ is $S$-perfect.
Thus we may apply Lemma \ref{lemma-derived-pushforward-rel-perfect}
to $K_n^\vee \otimes_{\mathcal{O}_X} E$ and we obtain an inverse system
$$
\ldots \to M_3 \to M_2 \to M_1
$$
of perfect complexes on $S$ with
$$
M_n = Rf_*(K_n^\vee \otimes_{\mathcal{O}_X}^\mathbf{L} E) =
Rf_*R\SheafHom(K_n, E)
$$
Moreover, the formation of these complexes commutes with any
base change, namely $Lg^*M_n =
Rf'_*((L(g')^*K_n)^\vee \otimes_{\mathcal{O}_{X'}}^\mathbf{L} L(g')^*E) =
Rf'_*R\SheafHom(L(g')^*K_n, L(g')^*E)$.

\medskip\noindent
As $K_n \to K$ induces an isomorphism on $\tau_{\geq -n}$, we see that
$K_n \to K_{n + 1}$ induces an isomorphism on $\tau_{\geq -n}$.
It follows that $K_{n + 1}^\vee \to K_n^\vee$
induces an isomorphism on $\tau_{\leq n}$ as
$K_n^\vee = R\SheafHom(K_n, \mathcal{O}_X)$.
Suppose that $E$ has tor amplitude in $[a, b]$ as a complex
of $f^{-1}\mathcal{O}_Y$-modules. Then the same is true after
any base change, see Lemma \ref{lemma-tor-independence-and-tor-amplitude}.
We find that
$K_{n + 1}^\vee \otimes_{\mathcal{O}_X} E \to
K_n^\vee \otimes_{\mathcal{O}_X} E$
induces an isomorphism on $\tau_{\leq n + a}$
and the same is true after any base change.
Applying the right derived functor $Rf_*$
we conclude the maps $M_{n + 1} \to M_n$
induce isomorphisms on $\tau_{\leq n + a}$
and the same is true after any base change.
Choose a distinguished triangle
$$
M_{n + 1} \to M_n \to C_n \to M_{n + 1}[1]
$$
Take $S'$ equal to the spectrum of the residue field at a point
$s \in S$ and pull back to see that
$C_n \otimes_{\mathcal{O}_S}^\mathbf{L} \kappa(s)$
has nonzero cohomology only in degrees $\geq n + a$. By
More on Algebra, Lemma
\ref{more-algebra-lemma-lift-perfect-from-residue-field}
we see that the perfect complex $C_n$ has tor amplitude in
$[n + a, m_n]$ for some integer $m_n$.
In particular, the dual perfect complex $C_n^\vee$ has tor amplitude in
$[-m_n, -n - a]$.

\medskip\noindent
Let $L_n = M_n^\vee$ be the dual perfect complex. The
conclusion from the discussion in the previous paragraph is that
$L_n \to L_{n + 1}$ induces isomorphisms on $\tau_{\geq -n - a}$.
Thus $L = \text{hocolim} L_n$ is pseudo-coherent, see
Lemma \ref{lemma-pseudo-coherent-hocolim}.
Since we have
$$
R\SheafHom(K, E) = R\SheafHom(\text{hocolim} K_n, E) =
R\lim R\SheafHom(K_n, E) = R\lim K_n^\vee \otimes_{\mathcal{O}_X} E
$$
(Cohomology, Lemma \ref{cohomology-lemma-colim-and-lim-of-duals})
and since $R\lim$ commutes with $Rf_*$ we find that
$$
Rf_*R\SheafHom(K, E) = R\lim M_n = R\lim R\SheafHom(L_n, \mathcal{O}_S) =
R\SheafHom(L, \mathcal{O}_S)
$$
This proves the formula over $S$. Since the construction of $M_n$ is
compatible with base chance, the formula continues to hold after
any base change.
\end{proof}

\end{document}
